\section{Details of the semidefinite representability arguments}
\label{repr:appendix}

\subsection{The real closed field inputs and their use}
\label{repr:field-inputs}

We record the external results needed in
Theorems~\ref{repr:threshold} and~\ref{repr:apex}. This also distinguishes the
copositive nonrepresentability theorem from the additional transfer to the
box.

A field is \emph{real closed} if it is ordered and every positive element has
a square root and every polynomial of odd degree has a root. The transfer
principle for real closed fields says that a first-order formula with real
parameters holds over $\R$ if and only if it holds over any real closed
extension of $\R$. In particular, copositivity, box nonnegativity of a
quadratic, and the equivalence of a semialgebraic set with a fixed finite
semidefinite lift transfer to such an extension. Positive semidefiniteness
can be expressed by the nonnegativity of finitely many principal minors.

For a real closed field $F$ and a divisible ordered abelian group $\Gamma$,
the Hahn field $F((\varepsilon^\Gamma))$ consists of formal series
$\sum_{a\in\Gamma}c_a\varepsilon^a$ with well-ordered support. The sign of a
nonzero series is the sign of its coefficient at the least exponent. This
field is real closed \citep[Section~2]{BodirskyKummerThom2024}.
Its group algebra $F[\Gamma]$ consists of the finite
such sums, regarded without reference to the ordering.

Call a function $f:\Gamma\to F$ \emph{positive definite} here if, for every
list of distinct $a_1,\ldots,a_r\in\Gamma$, the matrix
$(f(a_i+a_j))_{i,j}$ is positive definite. The use of sums of indices, rather
than differences, follows the ordered-field setting of
\citet{BodirskyKummerThom2024}.

\begin{lemma}[External inputs from Bodirsky--Kummer--Thom]
\label{repr:bkt-input}
The following statements hold.
\begin{enumerate}[label=(\roman*)]
\item For the Horn form $h(v)=v^\top Hv$, the polynomial
      $p(z)=h(z_1^2,\ldots,z_5^2)$ is nonnegative on $\R^5$, but
      $p(z_1^d,\ldots,z_5^d)$ is not a sum of squares for any integer $d\geq1$.
      Consequently $p$ is not a sum of squares in $\R[\Q^5]$.
\item There are a real closed extension $F\supseteq\R$ and an $F$-linear
      functional $T:F[\Q^5]\to F$ such that $T(a^2)>0$ for every nonzero
      $a\in F[\Q^5]$ and $T(p)<0$.
\item If $\Gamma$ is divisible and ordered, $F$ is real closed, and
      $f:\Gamma\to F$ is positive definite, the map
      $L_f(\sum_a c_a\varepsilon^a)=\sum_a f(a)c_a\varepsilon^a$ on
      $F((\varepsilon^\Gamma))$ is $F$-linear and completely positive:
      entrywise application preserves positive semidefinite matrices of
      every size. If $f(0)=1$, this map is unital.
\end{enumerate}
Statement (i) uses Example~3.4 and Remark~3.2, statement (ii) is the
positive-square separation in Theorem~3.7 as applied in Proposition~3.11,
and statement (iii) is Theorem~2.13(1) of
\citet{BodirskyKummerThom2024}.
\end{lemma}

Only these external inputs, transfer for real closed fields, and closed-cone
duality are used in the cube obstruction. Here is the passage from them to
\eqref{repr:negative-horn}. Normalize $T$ by its positive value $T(1)$ and set
\[
 f(a)=\frac{T(z^a)}{T(1)},\qquad a\in\Q^5.
\]
Then $f(0)=1$. For distinct $a_i$ and a nonzero vector $(c_i)\in F^r$,
\[
 \sum_{i,j}c_ic_j f(a_i+a_j)
 =\frac{T((\sum_i c_i z^{a_i})^2)}{T(1)}>0,
\]
so $f$ is positive definite. Since
$p=\sum_{i,j}H_{ij}z^{2e_i+2e_j}$, its negative value under $T$ gives
\eqref{repr:negative-horn}. The order on $\Q^5$ does not affect this
construction of $f$, so it may be chosen lexicographically afterward.
The coefficient map of part (iii) is then the map used in
\eqref{repr:coefficient-map}.

For clarity, the preservation of a semidefinite lift under this map has a
short direct proof. Suppose a semialgebraic set $S\subseteq\R^d$ has the
lift \eqref{eq:set:shadow}, and let $E$ be a real closed extension. Transfer
says that the same existential matrix formula describes $S(E)$. If
$(y,z)$ is a feasible witness over $E$, a unital real-linear completely
positive map $L:E\to E$ gives
\[
 0\preceq L\left(A_0+\sum_i y_iA_i+\sum_j z_jB_j\right)
 =A_0+\sum_i L(y_i)A_i+\sum_j L(z_j)B_j.
\]
Thus the coordinatewise image $L(y)$ belongs to $S(E)$. Unitality fixes the
constant matrix, real-linearity fixes its real coefficients, and complete
positivity preserves the matrix inequality. No order preservation of
multivariate polynomial evaluation is being assumed: the negative
evaluation in \eqref{repr:evaluation} is computed after mapping the
coefficients.

We also give an elementary verification that $H$ is copositive. For $v\geq0$,
\[
 v^\top Hv=\left(\sum_i v_i\right)^2
          -4\sum_{ij\in E(C_5)}v_iv_j.
\]
Fix the sum of the coordinates and maximize the edge sum on this compact
simplex. Among its maximizers choose one with the smallest support. If two
nonadjacent vertices have positive weight, move their combined weight to
the one with the larger weighted neighbor sum. Their mutual product does
not occur, so this does not reduce the edge sum, and it reduces the support.
This is a contradiction. The support of a minimal maximizer is therefore
a clique, which has at most two vertices in $C_5$. For such a support the
edge sum is at most one quarter of the squared coordinate sum. Hence
$v^\top Hv\geq0$. This verifies the copositivity input without any numerical
test.

\subsection{The simple-ray geometric condition}
\label{repr:simple-ray-proof}

Let $K\subseteq\R^m$ be full dimensional and pointed, and let $r$ span a
simple extreme ray. Write $\ell_2,\ldots,\ell_m$ for inner facet forms of
the $m-1$ incident facets. They are linearly independent: their common
kernel is the linear span of the ray, since intersecting the incident
facets recovers the ray. Choose $\ell_1$ strictly positive on
$K\setminus\{0\}$. Such a form exists for a pointed full-dimensional
polyhedral cone. It is not in the span of the incident facet forms because
it does not vanish on $r$. Thus
\[
 A:z\longmapsto(\ell_1(z),\ldots,\ell_m(z))
\]
is invertible. Normalize $r$ so that $\ell_1(r)=1$, and put
$v_k=A^{-1}e_k$ for $k=2,\ldots,m$.

Every nonincident facet form $g_j$ satisfies $g_j(r)>0$. Choose
$\delta>0$ such that, for every such facet,
\begin{equation}\label{repr:facet-bound}
 \delta\sum_{k=2}^5\abs{g_j(v_k)}\leq \tfrac12g_j(r).
\end{equation}
If there are no nonincident facets, any positive $\delta$ works. For a vector
$u\in\R_+^5$ with $u_k\leq\delta u_1$, first suppose $u_1>0$ and set
\[
 z=u_1r+\sum_{k=2}^5u_kv_k.
\]
Its incident facet values are $u_2,\ldots,u_5,0,\ldots,0$, all nonnegative.
For a nonincident facet, \eqref{repr:facet-bound} gives
\[
 g_j(z)\geq u_1g_j(r)-\delta u_1\sum_{k=2}^5\abs{g_j(v_k)}
           \geq\tfrac12u_1g_j(r)>0.
\]
Thus $z\in K$, and $(\ell_1,\ldots,\ell_5)(z)=u$. When $u_1=0$, all of
$u$ is zero and $z=0$ suffices. Each $\ell_i$ is nonnegative on $K$, so this
proves both inclusions in \eqref{repr:apex-condition}.

For the polytope consequence, use intrinsic affine coordinates if the
polytope is not full dimensional in its given ambient space. The cone
$K=\cone(\{1\}\times X)$ is closed and pointed, with dimension
$\dim X+1$. Its facets incident to the ray through $(1,v)$ correspond
exactly to facets of $X$ incident to $v$. The map from symmetric forms on
$K$ to quadratic polynomials on $X$, obtained by setting the first
coordinate to $1$, is a linear isomorphism. This proves the correspondence
used in Corollary~\ref{repr:simple-vertex}.

\subsection{The zero sets that describe proper faces}
\label{repr:face-proof}

Let $p\geq0$ be a nonzero quadratic on $C_4$, and let $F$ be a face of the
cube whose relative interior contains a zero $a$ of $p$. In affine
coordinates on $F$, local nonnegativity at $a$ implies zero gradient and a
positive semidefinite quadratic part $A_F$. Since $p$ is quadratic,
\[
 p(a+v)=v^\top A_Fv
\]
holds throughout the affine hull of $F$. Its zero set on $F$ is therefore
\[
 Z_F=F\cap(a+\ker A_F),
\]
a polytope. Every cube zero lies in the relative interior of a unique cube
face, so the whole zero set is the union of the finitely many $Z_F$ obtained
in this manner. Each has dimension at most three. For boundary faces this
follows from their dimension. For the full cube, $A_F$ cannot be zero,
because zero gradient and zero value at $a$ would then imply $p=0$.

To verify that the associated face of $\mathcal Q_4$ is precisely the
convex hull of the moments over this zero set, use a finite convex
representation by cube atoms. A nonnegative sum of their $p$-values is
zero exactly when every positive-weight atom has $p$-value zero. Thus
\[
 \{(m,Y)\in\mathcal Q_4:\inner{p}{(1,m,Y)}=0\}
 =\conv\{(x,xx^\top):x\in\textstyle\bigcup_F Z_F\}.
\]
For a polytope of affine dimension at most three, triangulation and
$\CP_r=\DNN_r$ for $r\leq4$ give a finite moment lift. An affine change of
coordinates $x=a+Bz$ acts linearly on the homogenized moments
$(1,z,zz^\top)$, so this conclusion applies to each $Z_F$ in its original
coordinates.

If compact sets $S_1,\ldots,S_r$ have finite lifts, Lemma~\ref{repr:normalization}
with $D=\{0\}$ gives finite lifts of their closed homogenized cones. The
system
\[
 (t_i,y_i)\in\widehat{S_i},\qquad
 \sum_i t_i=1,\qquad y=\sum_i y_i
\]
is a finite lift of $\conv(\bigcup_i S_i)$. At zero weight its visible
contribution is zero; at positive weight it contributes a weighted point
of $S_i$. This justifies the convex-hull step even when the original lifts
have unbounded auxiliary variables.

Finally, let $F$ be an arbitrary proper face of a compact convex set $S$ in
its affine hull, and take $a\in\operatorname{relint}F$. A proper supporting
hyperplane of $S$ through $a$ cuts out an exposed proper face $E$ containing
$F$. The inclusion follows because a linear function attaining its minimum
at a relative interior point of $F$ is constant on $F$. If $F\neq E$, it is
a proper face of $E$, so repeat the argument in the affine hull of $E$.
At every step the containing face decreases in dimension. Eventually it
equals $F$. After the first step, all these affine sections preserve the
finite lift already established for $E$. This proves
Proposition~\ref{repr:proper-faces} for nonexposed faces as well.
